% Part: incompleteness
% Chapter: incompleteness-provability
% Section: provability-conditions

\documentclass[../../../include/open-logic-section]{subfiles}

\begin{document}

\olfileid[id]{inc}{inp}{prc}

\olsection{Syarat \usetoken{S}{derivability} bagi $\Th{PA}$}

Aritmetika Peano, atau $\Th{PA}$, adalah teori yang memperluas $\Th{Q}$
dengan aksioma induksi bagi semua !!{formula}s. Dengan kata lain, pada
$\Th{Q}$ ditambahkan aksioma-aksioma berbentuk
\[
(!A(0) \land \lforall[x][(!A(x) \lif !A(x'))]) \lif \lforall[x][!A(x)]
\]
untuk setiap !!{formula}~$!A$. Perhatikan bahwa ini sebenarnya suatu
\emph{skema}, yakni tak berhingga banyaknya aksioma (dan ternyata $\Th{PA}$
{\em tidak} dapat diaksiomatisasi secara berhingga). Namun, karena kita
dapat menentukan secara efektif apakah suatu untai simbol merupakan contoh
aksioma induksi, himpunan aksioma bagi $\Th{PA}$ dapat dihitung. $\Th{PA}$
adalah teori yang jauh lebih kuat daripada~$\Th{Q}$. Misalnya, dengan
menggunakan induksi dengan cara biasa, kita dapat dengan mudah membuktikan
bahwa penjumlahan dan perkalian bersifat komutatif. Bahkan, sebagian besar
argumen finitistik dalam teori bilangan dan kombinatorika dapat dilakukan
dalam~$\Th{PA}$.

Karena $\Th{PA}$ diaksiomatisasi secara komputabel, predikat
!!{derivability} $\Prf[\Th{PA}](x,y)$ dapat dihitung dan karenanya
direpresentasikan dalam~$\Th{Q}$ (dan dengan demikian dalam~$\Th{PA}$).
Seperti sebelumnya, kita akan menggunakan $\OPrf[\Th{PA}](x,y)$ untuk
menyatakan formula yang merepresentasikan relasi tersebut. Misalkan
$\OProv[\Th{PA}](y)$ adalah formula
$\lexists[x][\OPrf[\Th{PA}](x,y)]$, yang secara intuitif mengatakan,
``$y$~!!{derivable} dari aksioma-aksioma $\Th{PA}$.'' Alasan kita
memerlukan sedikit lebih banyak daripada aksioma-aksioma $\Th{Q}$ ialah
bahwa kita perlu mengetahui bahwa teori yang digunakan cukup kuat untuk
!!{derive} beberapa fakta dasar tentang predikat !!{derivability} ini.
Fakta-fakta yang kita perlukan adalah sebagai berikut:
\begin{enumerate}
\item[P1.] Jika $\Th{PA} \Proves !A$, maka $\Th{PA} \Proves
  \OProv[\Th{PA}](\gn{!A})$.
\item[P2.] Untuk semua !!{formula}s $!A$ dan $!B$,
  \[
  \Th{PA} \Proves \OProv[\Th{PA}](\gn{!A \lif !B}) \lif
  (\OProv[\Th{PA}](\gn{!A}) \lif \OProv[\Th{PA}](\gn{!B})).
  \]
\item[P3.] Untuk setiap !!{formula}~$!A$,
  \[
  \Th{PA} \Proves \OProv[\Th{PA}](\gn{!A})
  \lif \OProv[\Th{PA}](\gn{\OProv[\Th{PA}](\gn{!A})}).
  \]
\end{enumerate}
Satu-satunya cara untuk memverifikasi bahwa ketiga sifat ini berlaku ialah
menguraikan !!{formula} $\OProv[\Th{PA}](y)$ dengan cermat dan menggunakan
aksioma-aksioma $\Th{PA}$ untuk menguraikan !!{derivation}s formal yang
relevan. Syarat~(1) dan~(2) mudah; sesungguhnya syarat~(3)-lah yang
memerlukan usaha. (Pikirkan jenis usaha apa yang diperlukannya \dots)
Menjabarkan perinciannya akan menjemukan dan tidak menarik, sehingga di sini
kita meminta Anda menerima bahwa $\Th{PA}$ memiliki ketiga sifat di atas.
Pilihan yang wajar bagi $\OProv[\Th{PA}](y)$ juga akan memenuhi
\begin{enumerate}
\item[P4.] Jika $\Th{PA} \Proves \OProv[\Th{PA}](\gn{!A})$, maka
  $\Th{PA} \Proves !A$.
\end{enumerate}
Namun, kita tidak akan memerlukan fakta ini.

\begin{digress}
Kebetulan, G\"odel juga bermalas-malasan seperti kita sekarang. Pada akhir
makalah tahun 1931 itu, ia membuat sketsa pembuktian teorema ketaklengkapan
kedua dan menjanjikan perinciannya dalam makalah berikutnya. Ia tidak pernah
sempat melakukannya; karena setiap orang yang memahami argumennya percaya
bahwa pembuktian itu dapat dijabarkan, ia tidak perlu mengisi perinciannya.
\end{digress}

\end{document}
